\section{Counting Numbers}
\label{sec:cses-counting-numbers}

\problemstatement{Count the integers in the range $[a,b]$ that do not
contain two adjacent equal digits (e.g.\ $33$ and $121\underline{22}$ are
excluded). Here $0\le a\le b\le10^{18}$.}

\begin{patternbox}
\emph{``Count numbers up to $N$ with a digit property''} is
\textbf{digit DP}. Reduce $[a,b]$ to $f(b)-f(a-1)$ where $f(x)$ counts
valid numbers in $[0,x]$, then build $x$'s digits left to right, tracking
just enough state to enforce ``no two adjacent equal digits.''
\end{patternbox}

\subsection*{State for digit DP}

Write $x$ as a digit string and fill positions left to right. The state
is:
\begin{itemize}
  \item \textbf{position} in the digit string,
  \item \textbf{previous digit} placed (to forbid an equal neighbour),
  \item \textbf{tight}: whether the prefix so far equals $x$'s prefix (if
        so, the next digit is capped by $x$'s digit),
  \item \textbf{started}: whether we have placed a nonzero digit yet (to
        handle leading zeros, which must not count as ``adjacent equal'').
\end{itemize}
At each position we try every allowed next digit $d$ (up to the cap if
tight), skipping $d$ equal to the previous digit once the number has
started. Memoise over (position, previous, started) for the non-tight
states.

\begin{ideabox}[Tight and Started Are the Two Universal Digit-DP Flags]
\textbf{Tight} keeps the enumeration within $[0,x]$; it is true only along
the single prefix that hugs $x$. \textbf{Started} distinguishes real
leading digits from padding zeros, so a leading zero does not falsely
collide with the next digit. Almost every digit-DP carries these two.
\end{ideabox}

\subsection*{Algorithm}

\begin{enumerate}
  \item Define $f(x)$ = count of valid numbers in $[0,x]$ by digit DP
        over $x$'s digits with state (pos, prev, tight, started).
  \item At each position try digits $0\ldots(\text{cap})$; skip $d=$
        prev when started; recurse with updated flags.
  \item Answer $f(b)-f(a-1)$ (with $f(-1)=0$ when $a=0$).
\end{enumerate}

\subsection*{Dry run}

\dryrun{Count valid numbers in $[0,23]$.}

\begin{itemize}
  \item Invalid two-digit numbers $\le23$ with equal adjacent digits:
        $11,22$.
  \item All $24$ numbers $0\ldots23$ minus the $2$ invalid ones $=22$
        valid numbers.
  \item The digit DP arrives at the same $22$ by never placing a digit
        equal to its left neighbour.
\end{itemize}

\subsection*{C++ implementation}

\begin{lstlisting}
#include <bits/stdc++.h>
using namespace std;

string s;
long long memo[20][11][2];      // [pos][prev+1][started], for tight=false

long long solve(int pos, int prev, bool tight, bool started) {
    if (pos == (int)s.size()) return 1;          // one valid number formed
    if (!tight && memo[pos][prev + 1][started] != -1)
        return memo[pos][prev + 1][started];

    int cap = tight ? (s[pos] - '0') : 9;
    long long total = 0;
    for (int d = 0; d <= cap; d++) {
        if (started && d == prev) continue;      // no equal adjacent digits
        bool nstarted = started || (d > 0);
        int nprev = nstarted ? d : -1;           // leading zeros carry prev=-1
        total += solve(pos + 1, nprev, tight && (d == cap), nstarted);
    }
    if (!tight) memo[pos][prev + 1][started] = total;
    return total;
}

long long f(long long x) {
    if (x < 0) return 0;
    s = to_string(x);
    memset(memo, -1, sizeof(memo));
    return solve(0, -1, true, false);
}

int main() {
    long long a, b;
    cin >> a >> b;
    cout << f(b) - f(a - 1) << "\n";
    return 0;
}
\end{lstlisting}

\begin{complexitybox}
\textbf{Time Complexity.} $O(\text{digits}\times10\times10)$ per call --
about $20\cdot10\cdot10$ states, each trying $10$ digits: constant for
64-bit inputs. \textbf{Space Complexity.} $O(\text{digits}\times10\times
2)$ for the memo.
\end{complexitybox}

\begin{takeawaybox}
Digit DP counts numbers with a property by building them digit by digit
under the two universal flags \emph{tight} and \emph{started}, memoising
the free (non-tight) states. Range queries become $f(b)-f(a-1)$. This
template adapts to any per-digit constraint by changing the ``skip this
digit'' rule.
\end{takeawaybox}
